\section{ARCP: Acquisition-Robust Camouflage Perturbations}
\label{sec:defense}

ARCP-v2 is designed for the revised threat model rather than inferred from the
legacy 40-identity results. A uniform high-frequency perturbation is vulnerable
to transcoding and resizing, a uniform amplitude wastes visibility budget in
smooth facial regions, and a T0-only objective ignores both acquisition and the
reader population. ARCP-v2 addresses these mechanisms with band limitation,
texture allocation, and an acquisition-averaged reader-diverse objective. These
are design choices, not evidence of superiority; Section~\ref{sec:results}
withholds that claim until the paired 90-identity evaluation finishes.

\paragraph{Band-limited parameterization with perceptual amplitude allocation.}
The perturbation is parameterized in a low-resolution latent grid $z_i \in [-1,1]^{64\times64\times3}$ and rendered by bilinear upsampling, $\delta_i = A(x_i) \odot U_8(z_i)$, which hard-limits the perturbation spectrum to frequencies that survive resampling, JPEG quantization, and super-resolution by construction. The amplitude map
\begin{equation}
A(x) \;=\; \epsilon_{\min} + (\epsilon_{\max}-\epsilon_{\min})\,\hat{t}(x),
\label{eq:jnd}
\end{equation}
allocates budget by texture: $\hat{t}(x)$ is the Sobel-energy map of $x$, Gaussian-smoothed and normalized at its 95th percentile, so that $\epsilon_{\min}=2/255$ is spent on smooth regions and up to $\epsilon_{\max}=16/255$ inside textured regions where contrast masking hides it. The allocation is fixed before optimization and costs one forward pass; no perceptual network is invoked.

\paragraph{Expectation over the acquisition distribution.}
Each ARCP-v2 optimization step samples one differentiable proxy from
$\mathcal{T}$: diff-JPEG at $q\sim U[65,95]$, bicubic down-and-up resampling at
$s\sim U[0.35,1.0]$, a half-resolution super-resolution surrogate, or an
affine face crop at $f\sim U[0.7,1.0]$. These directly approximate T1--T4.
T5 and T6 are covered only approximately in latent space: with probability one
half the denoising loss uses $t\sim U[300,500]$, rather than differentiating
through the exact img2img or purification pipeline. Their robustness is
therefore an empirical question, not guaranteed by the training objective.

\paragraph{Reader-diverse objective.}
ARCP optimizes
\begin{equation}
\max_{z}\;\; \mathbb{E}_{T}\big[\mathcal{L}_{\mathrm{den}}(T(\tilde{x}))\big] \;-\; \lambda\,\mathbb{E}_{T}\big[\mathrm{relu}\big(\cos\big(f_{\mathrm{face}}(T(\tilde{x})),\,\bar{e}_s\big)-m\big)\big],
\label{eq:arcp}
\end{equation}
where expectations are over proxy compositions $T\sim\mathcal{T}$,
$\mathcal{L}_{\mathrm{den}}$ is the denoising-MSE poisoning objective,
$f_{\mathrm{face}}$ is an InceptionResnetV1 face-metric surrogate, and
$\bar{e}_s$ is the clean-gallery centroid. The frozen revised setting uses
$m=.35$, $\lambda=.75$, and 200 optimization steps. Evaluation readers are not
optimized against. The first term targets training-based customization while
the second explicitly repels the identity representation consumed by zero-shot
readers.

\paragraph{Scope.}
ARCP-v2 does not claim in advance that an invisible perturbation survives every
operation. T0 is the unaltered baseline and each of T1--T6 independently consumes
the same protected source. Its testable objective is a nonzero paired protection
gap under each intervention while preserving image quality; the completed run,
not the optimization objective, determines whether that objective is met.
